\chapter{Problem statement}
\label{ch:problem_statement}

This chapter formalizes the search problem studied in this thesis, introduces the notation used throughout, and states the research goal in a precise form. The focus is on a minimal but fully specified two-dimensional model of random search with finite detection radius, in which the effect of the Lévy tail exponent $\alpha$ can be benchmarked under controlled assumptions.

%-----------------------------------------------------------
\section{Search domain, targets, and trajectory}
The searcher operates in the Euclidean plane~$\mathbb{R}^2$.
In principle, the domain may be bounded (a finite region~$\Omega$ with
periodic or reflecting boundary conditions) or unbounded.  The
implementation used in this thesis (Chapter~\ref{ch:methodology})
adopts an effectively infinite plane realized via on-demand
generation of square chunks, thereby avoiding the finite-size
and boundary effects associated with periodic domains.

Targets are distributed as a homogeneous Poisson point process
with intensity~$\rho$ (targets per unit area).  The probability
of finding exactly~$k$ targets in a region of area~$S$ is
$P(k) = (\rho S)^k e^{-\rho S}/k!$.  The expected number of
targets in any region is $\rho S$, and target positions are
independent and uniformly distributed.

A searcher moves at constant unit speed and produces a
trajectory $x(t)\in\mathbb{R}^2$ for $t\ge 0$, so that
distance and time are interchangeable.
In simulations, each flight is decomposed into segments of
length~$\Delta x$ (the scan resolution), and detections are
checked at each segment.  The total search budget is specified
as a number of flights~$N_{\mathrm{flights}}$; together with the
exponent-dependent mean flight length, this determines the total
distance traveled.

\begin{Def}[Detection event]
Fix a detection radius $r_v>0$. A target $y_i$ is detected at time $t$ if
\begin{equation}
    \|x(t)-y_i\|\le r_v .
    \label{eq:detection_condition}
\end{equation}
\end{Def}

We emphasize that the finite radius $r_v$ is essential in 2D: if targets were modeled as mathematical points and detections required exact hits, then a continuous trajectory would detect a target with probability zero. The model therefore treats targets as detectable discs of radius $r_v$ (or, equivalently, point targets with an observation radius~$r_v$).

%-----------------------------------------------------------
\section{Motion model: incremental relocations}
Between detections, the searcher executes \emph{relocations}. Each relocation begins by sampling a direction and a relocation length. Let $\theta_k \sim \mathrm{Unif}(0,2\pi)$ be an independent heading angle, and let $\ell_k>0$ be the relocation length sampled from a parameterized family of heavy-tailed laws. The searcher then moves toward the point
\begin{equation}
    x_k^{\star} = x(t_k) + \ell_k \, (\cos\theta_k,\sin\theta_k),
    \label{eq:relocation_target_point}
\end{equation}
but the motion is \emph{incremental}: the flight is decomposed into segments of length~$\Delta x$, and the detection condition~\eqref{eq:detection_condition} is checked along each segment.  Thus relocations can be \emph{interrupted} by detections before the planned length $\ell_k$ is completed.  (In subsequent chapters, we use ``flight'' as a synonym for ``relocation.'')

\begin{Def}[Heavy-tailed L\'evy relocation family]
Fix $\alpha\in(0,2)$. A relocation distribution is called L\'evy-type if it has a power-law tail
\begin{equation}
    \Pr(\ell > s)\sim C\, s^{-\alpha}\qquad \text{as } s\to\infty,
    \label{eq:levy_tail}
\end{equation}
for some constant $C>0$. In practice, simulations use a truncated version on $[l_{\min}, l_{\max}]$ to ensure a finite mean flight length and well-defined sampling.
\end{Def}

For comparison baselines, the same incremental motion rule is used with alternative choices of relocation lengths:
\begin{itemize}
    \item \textbf{Brownian-like baseline:} short-tailed relocation lengths (e.g., exponential or Gaussian);
    \item \textbf{Ballistic baseline:} very long relocations (effectively $\alpha\to 0$ or $\mu\to 1$ in power-law notation), producing long straight runs.
\end{itemize}

%-----------------------------------------------------------
\section{Target interaction rules}
The model distinguishes two interaction regimes.

\begin{Def}[Destructive targets]
In destructive search, each target can be detected at most once. After the first detection of $y_i$, the target is removed from the active set for the remainder of the simulation.
\end{Def}

\begin{Def}[Nondestructive targets with explicit post-hit rule]
In nondestructive search, targets may be revisited. To avoid an ill-posed “stay forever” behavior, a post-hit departure mechanism is specified explicitly. Two common choices are:
\begin{enumerate}
    \item \textbf{Depletion/renewal time} $\tau \ge 0$: after a detection, the target is inactive for time $\tau$ and becomes detectable again afterwards.
    \item \textbf{Jump-off distance} $l_c>0$: after a detection, the searcher is repositioned to a point at distance $l_c$ from the target (in a random direction), ensuring departure. It is convenient to use the dimensionless parameter
    \begin{equation}
        \varepsilon := \frac{l_c}{r_v}-1 .
        \label{eq:epsilon_jumpoff}
    \end{equation}
\end{enumerate}
\end{Def}

These interaction rules are treated as first-class model parameters
rather than hidden implementation details, because they can change
qualitative conclusions about optimality.  In the implementation
(Chapter~\ref{ch:methodology}), the teleport strategy uses
$\varepsilon = 4$ ($l_c = 5\,r_v$), while the non-destructive
Viswanathan restart corresponds to the singular limit
$\varepsilon \to -1$ ($l_c \to 0$); in this limit, a minimal
``just-visited'' flag prevents the forager from re-detecting
the same target on the immediately following flight step
(see Chapter~\ref{ch:methodology} for details).

%-----------------------------------------------------------
\section{Performance criteria}
Let $\{t_j\}_{j\ge 1}$ be the (random) sequence of detection times during a run.

\begin{Def}[Mean first-passage time (MFPT)]
The mean first-passage time is the expected total distance
traveled before the first detection (equivalent to time because
the searcher moves at constant unit speed):
\begin{equation}
    \mathrm{MFPT} := \mathbb{E}[D_{\mathrm{first}}],
    \label{eq:mfpt}
\end{equation}
where $D_{\mathrm{first}}$ is the path length from the start of a
run to the first detection, and the expectation is taken over random
relocations and target placements.  The acronym MFPT follows
standard usage~\cite{viswanathan1999optimizing}.
\end{Def}

\begin{Def}[Unique-capture rate]
Over $R$ independent runs, the \emph{pooled}
unique-capture rate is defined as
\begin{equation}
    \eta := \frac{\sum_{r=1}^{R} N_{\mathrm{unique}}^{(r)}}
                 {\sum_{r=1}^{R} D_{\mathrm{total}}^{(r)}},
    \label{eq:eta_formal}
\end{equation}
where $N_{\mathrm{unique}}^{(r)}$ is the number of distinct
targets found in run~$r$ and $D_{\mathrm{total}}^{(r)}$ is the
total distance traveled.  This pooled estimator is preferred over
the per-run mean because it is robust to heavy-tailed runs
(see Chapter~\ref{ch:methodology}).  The dimensionless form is
$\lambda\eta$, where $\lambda = 1/(2\rho r_v)$ is the mean free
path.
\end{Def}

In addition to mean performance, robustness is quantified by the
dispersion of outcomes across independent runs and by convergence
studies with respect to the number of flights and runs.

%-----------------------------------------------------------
\section{Formal goal of the study}
The study aims to build a transparent numerical benchmark and, within it, to identify which Lévy exponent is optimal across regimes.

\begin{Def}[Parameter space]
Let $\Theta$ denote the set of environmental, numerical, and
interaction parameters:
\begin{equation}
    \Theta := \bigl\{\rho,\; r_v,\; l_{\min},\; l_{\max},\;
    \Delta x,\; N_{\mathrm{flights}},\;
    \mathrm{strategy}\bigr\}.
    \label{eq:parameter_space}
\end{equation}
Here \emph{strategy} takes one of four values: no-interrupt,
teleport, Viswanathan destructive, or Viswanathan non-destructive.
The key dimensionless control parameter is the sparsity ratio
$\lambda/r_v = 1/(2\rho r_v^2)$.  Note that the no-interrupt
strategy is inherently destructive (each target is counted at
most once).
\end{Def}

\begin{Def}[Objective-dependent optimal exponent]
Fix a choice of objective $\mathcal{J}\in\{\mathrm{MFPT},\, \eta\}$
and fix parameters $\theta\in\Theta$. The objective-dependent
optimal exponent is
\begin{equation}
    \alpha^\star(\theta;\mathcal{J}) \in \arg\max_{\alpha\in(0,2)}
    \, \mathcal{J}(\alpha;\theta),
    \label{eq:alpha_star}
\end{equation}
where for MFPT the maximization is replaced by minimization.
Equivalently, in the $\mu$-convention
(Chapter~\ref{ch:methodology}), $\mu^* = \alpha^* + 1$.
\end{Def}

\paragraph{Overall purpose (formal statement).}
For each regime $\theta\in\Theta$, the purpose of this thesis
is to:
\begin{enumerate}
    \item implement a reproducible 2D simulator on an infinite
          plane with Poisson-distributed targets and four
          post-detection strategies;
    \item estimate $\mathrm{MFPT}(\mu;\theta)$ and
          $\lambda\eta(\mu;\theta)$ via Monte Carlo simulation;
    \item determine $\mu^*(\theta;\mathcal{J})$ for both
          objectives and assess convergence with respect to the
          number of flights and runs;
    \item map the dependence of $\mu^*$ on the sparsity ratio
          $\lambda/r_v$ across two orders of magnitude and on the
          truncation parameters $l_{\min}/r_v$ and
          $l_{\max}/\lambda$.
\end{enumerate}

\paragraph{Problem statement (one sentence).}
Given a two-dimensional random search model on an infinite Poisson
target field with finite detection radius and explicit
post-detection strategies, determine---by systematic Monte Carlo
simulation---how the optimal L\'evy exponent $\mu^*$ depends on
the sparsity ratio $\lambda/r_v$, the truncation parameters, and
the choice of strategy.
